\section{Natural language requirements}


CodeKataBattle (CKB) is a new platform that helps students improve their software development skills by training with peers on code kata1 . Educators use the platform to challenge students by creating code kata battles in which teams of students can compete against each other, thus proving (and improving) their skills. 


A code kata battle is essentially a programming exercise in a programming language of choice (e.g., Java, Python). The exercise includes a brief textual description and a software project with build automation scripts (e.g., a Gradle project in case of Java sources) that contains a set of test cases that the program must pass, but without the program implementation. Students are asked to complete the project with their code. In particular, groups of students participating in a battle are expected to follow a test-first approach2 and develop a solution that passes the required tests. Groups deliver their solution to the platform (by the end of the battle). At the end of the battle, the platform assigns scores to groups to create a competition rank. 


Each battle created by an educator belongs to a specific tournament. Tournaments are created by an educator who can then grant to other colleagues the permission to create battles within the context of a specific tournament. When a new tournament is created, all students subscribed to the CKB platform are notified and they can subscribe by a given deadline. If they subscribe, they are notified of all upcoming battles created within that tournament. To create a new battle, an educator uses the CKB platform to perform the following steps: 
\begin{itemize}
    \item upload the code kata (description and software project, including test cases and build automation scripts), 
    \item set minimum and maximum number of students per group, 
    \item set a registration deadline, 
    \item set a final submission deadline,
    \item set additional configurations for scoring (see further details in the following). 
\end{itemize}


After creation of a battle, students use the platform to form teams for that battle. In particular, each student can join a battle on his/her own or by inviting other students (respecting the minimum and maximum number of students per group set for that battle). When the registration deadline expires, the platform creates a GitHub repository containing the code kata and then sends the link to all students who are members of subscribed teams. At this point, students can start working on the project. Students are asked to fork the GitHub repository of the code kata and set up an automated workflow through GitHub Actions that informs the CKB platform (through proper API calls) as soon as students push a new commit into the main branch of their repository. Thus, each push before the deadline triggers the CKB platform, which pulls the latest sources, analyzes them, and runs the tests on the corresponding executables to calculate and update the battle score of the team. 


The score is a natural number between 0 and 100 determined by considering some mandatory factors evaluated in a fully automated way, and optional factors evaluated manually by educators. Mandatory automated evaluation includes: 
\begin{itemize}
    \item functional aspects, measured in terms of number of test cases that pass out of all test cases (the higher the better); 
    \item timeliness, measured in terms of time passed between the registration deadline and the last commit (the lower the better); 
    \item quality level of the sources, extracted through static analysis tools that consider multiple aspects such as security, reliability, and maintainability (the higher the better). Aspects are selected by the educator at battle creation time.
\end{itemize}
Optional manual evaluation includes: 
\begin{itemize}
    \item personal score assigned by the educator, who checks and evaluates the work done by students (the higher the better). 
\end{itemize}

The CKB platform automatically updates the battle score of a team as soon as new push actions on GitHub are performed. So, both students and educators involved in the battle can see the current rank evolving during the battle. When the submission deadline expires, there is a consolidation stage in which, if manual evaluation is required, the educator uses the CKB platform to go through the sources produced by each team to assign his/her score. Once the consolidation stage finishes, all students participating in the battle are notified when the final battle rank becomes available. 


At the end of each battle, the platform updates the personal tournament score of each student, that is, the sum of all battle scores received in that tournament. Thus, for each tournament, there is a rank that measures how a student's performance compares to other students in the context of that tournament. This information is available for all students and educators subscribed to the CKB platform, that is, all users can see the list of ongoing tournaments as well as the corresponding tournament rank. 


When an educator closes a tournament, as soon as the final tournament rank becomes available the CKB platform notifies all students involved in the tournament.


The CKB platform also includes gamification badges: elements in the form of rewards that represent the achievements of individual students. Badges are defined by educators when they create a tournament. Each badge has a title (e.g., “top committer”) and one or more rules (i.e., simple Boolean properties) that must be fulfilled to achieve the badge. Thus, each badge is assigned to one or more students depending on the rules checked at the end of the tournament. The following are examples of possible badges: 
\begin{itemize}
    \item Title: tournament participant 
o Rules: { tot\_attended\_battles > 0 } 
    \item Title: top committer 
o Rules: { tot\_commits\_student == max\_tot\_commits } 
\end{itemize}


Where \textit{tot\_attended\_battles}, \textit{tot\_commits\_student} and \textit{max\_tot\_commits} are pre-defined variables that represent, respectively, the total number of battles the student has been involved in, the number of commits carried out by the student and the maximum total number of commits considering all students. 


As mentioned above, educators can create new badges and define new rules as well as new variables associated with them. Variables can represent any piece of information available in CKB relevant for scoring. It is up to you to identify such information and a simple language for creating new variables, rules, and badges. 


Badges can be visualized by all users. In particular, both students and educators can see collected badges when they visualize the profile of a student. 
